\section{Alaee 表 1 的完整推导——长波长近似}
%=============================================================================

本章推导 Alaee 2018 \textbf{表 1}（长波长近似多极矩）。
先给出目标公式一览，再从多极积分核出发取 $kr'\ll1$ 近似，完整推出四条公式。

\subsection{表 1 公式一览}

\begin{table}[htbp]
\centering
\small
\begin{tabular}{ll}
\toprule
\textbf{多极} & \textbf{表达式（近似）} \\
\midrule
ED & $\displaystyle p_\alpha \approx -\frac{1}{i\omega}\Bigg\{\int d^3r\,J_\alpha
  + \frac{k^2}{10}\int d^3r\bigl[(\mathbf{r}\cdot\mathbf{J})r_\alpha - 2r^2 J_\alpha\bigr]\Bigg\}$ \\[1.2em]
MD & $\displaystyle m_\alpha \approx \frac12\int d^3r\,(\mathbf{r}\times\mathbf{J})_\alpha$ \\[1.2em]
EQ & $\displaystyle Q^e_{\alpha\beta} \approx -\frac{1}{i\omega}\Bigg\{
  \int d^3r\bigl[3(r_\beta J_\alpha + r_\alpha J_\beta) - 2(\mathbf{r}\cdot\mathbf{J})\delta_{\alpha\beta}\bigr]$ \\[0.5em]
  & $\displaystyle\qquad + \frac{k^2}{14}\int d^3r\bigl[4r_\alpha r_\beta(\mathbf{r}\cdot\mathbf{J})
  - 5r^2(r_\alpha J_\beta + r_\beta J_\alpha) + 2r^2(\mathbf{r}\cdot\mathbf{J})\delta_{\alpha\beta}\bigr]\Bigg\}$ \\[1.2em]
MQ & $\displaystyle Q^m_{\alpha\beta} \approx \int d^3r\bigl\{
  r_\alpha(\mathbf{r}\times\mathbf{J})_\beta + r_\beta(\mathbf{r}\times\mathbf{J})_\alpha\bigr\}$ \\[0.5em]
\bottomrule
\end{tabular}
\caption*{表 1：长波长近似多极矩（Alaee 2018 表 1），$\mathbf{J}\equiv\mathbf{J}^\omega$ 为时谐电流密度。}
\end{table}

\begin{noteworthy}
\textbf{结构观察}：表 1 的积分核全是\textbf{常数}（$1$、$\frac13$、$\frac1{15}$、$\frac1{105}$ 藏在系数里）。
这些常数正是球 Bessel 函数在 $kr\to0$ 时的首项——表 1 是表 2（$j_l$ 核）的 $kr\to0$ 极限。
\end{noteworthy}

\subsection{长波长近似的物理图像}

\begin{stepbox}{长波长近似是什么？}
散射体特征尺度为 $a$，入射波长为 $\lambda$。当 $ka \ll 1$（粒子远小于波长）时：
\begin{itemize}
  \item 粒子内部各点电流感受到的相位差 $\sim ka \ll 1$，可以忽略
  \item 因此 $j_l(kr') \approx (kr')^l/(2l+1)!!$ 只保留首项，积分核退化为纯多项式 $\mathbf{r}'^{\,l}$
  \item 多极矩退化为电流的坐标矩
\end{itemize}
\textbf{适用边界}：$ka\gtrsim0.5$ 时近似失效（Alaee Fig.1：电偶极、磁偶极误差可超 100\%）。
\end{stepbox}

\subsection{电偶极矩（表 1 ED）}

\subsubsection{主项：从 $\rho$ 积分到 $\mathbf{J}$ 积分}

电荷密度与电流的关系（连续性方程）：$\nabla\cdot\mathbf{J} = i\omega\rho$，即 $\rho = \dfrac{1}{i\omega}\nabla\cdot\mathbf{J}$。

偶极矩定义 $\mathbf{p} = \int \mathbf{r}\,\rho\,d^3r$，逐分量：
\begin{align}
  p_\alpha &= \int r_\alpha\,\rho\,d^3r
    = \frac{1}{i\omega}\int r_\alpha\,\nabla\cdot\mathbf{J}\,d^3r \notag\\
  &= -\frac{1}{i\omega}\int \mathbf{J}\cdot\nabla r_\alpha\,d^3r
    \qquad(\text{分部积分，表面项消失}) \notag\\
  &= -\frac{1}{i\omega}\int J_\alpha\,d^3r
\end{align}

\begin{stepbox}{为什么表面项为零？}
分部积分 $\int_V r_\alpha\nabla\cdot\mathbf{J}\,d^3r = \oint_{\partial V}r_\alpha\mathbf{J}\cdot d\mathbf{S} - \int_V\mathbf{J}\cdot\nabla r_\alpha\,d^3r$。
$\mathbf{J}$ 只在粒子内部非零，把积分体积 $V$ 取到粒子外足够远处，表面项为零。
\end{stepbox}

\textbf{于是表 1 的 ED 主项}：
\begin{equation}
  p_\alpha^{(\text{主})} = -\frac{1}{i\omega}\int J_\alpha\,d^3r
\end{equation}

\subsubsection{toroidal 修正项（表 1 ED 的 $k^2$ 项）}

表 1 的 ED 还包含 $\dfrac{k^2}{10}\int[(\mathbf{r}\cdot\mathbf{J})r_\alpha - 2r^2J_\alpha]$。
这正是\textbf{环形偶极矩（toroidal dipole）}的贡献。

\begin{stepbox}{toroidal 偶极矩}
标准定义（Radescu--Vaman，Alaee 引用的文献）：
\begin{equation}
  T_\alpha = \frac{1}{10c}\int\Bigl[(\mathbf{r}\cdot\mathbf{J})r_\alpha - 2r^2J_\alpha\Bigr]d^3r
  \label{eq:toroidal_def}
\end{equation}
它描述"环形电流绕一个轴打转"的模式——在长波长近似下与电偶极辐射\textbf{不可区分}，
其辐射场与原电偶极相干叠加，等效为修正的偶极矩：
\begin{equation}
  p_\alpha^{\text{eff}} = p_\alpha + ik\,T_\alpha
\end{equation}
\end{stepbox}

验证表 1 的系数：
\begin{equation}
  -\frac{1}{i\omega}\cdot\frac{k^2}{10}\int\bigl[(\mathbf{r}\cdot\mathbf{J})r_\alpha - 2r^2J_\alpha\bigr]
  = -\frac{1}{i\omega}\,k^2c\,T_\alpha
  = -\frac{\omega^2/c}{i\omega}T_\alpha
  = \frac{i\omega}{c}T_\alpha = ik\,T_\alpha
\end{equation}

\begin{nontrivial}
\textbf{重要}：表 1 的 ED 实际上是"电偶极 + 环形偶极"的\textbf{合并}形式——
它把 $p$ 和 $ikT$ 合成一个"有效偶极矩"。\\
长波长近似下 toroidal 贡献不可与偶极分离，必须一起给出。
\end{nontrivial}

\subsection{磁偶极矩（表 1 MD）}

磁偶极矩是电流环流的度量，标准定义：
\begin{equation}
  \mathbf{m} = \frac12\int \mathbf{r}\times\mathbf{J}\,d^3r
  \;\Longrightarrow\;
  m_\alpha = \frac12\int (\mathbf{r}\times\mathbf{J})_\alpha\,d^3r
\end{equation}

\begin{stepbox}{物理解释}
一段电流环 $I$ 围出面积 $\mathbf{S}$，则 $|\mathbf{m}| = IS$。
对于任意电流分布，$\frac12\int\mathbf{r}\times\mathbf{J}\,d^3r$ 正是把所有微小电流环叠加的结果。
\end{stepbox}

\subsection{电四极矩（表 1 EQ）}

电四极矩的定义（无迹对称）：
\begin{equation}
  Q^e_{\alpha\beta} = \int \bigl(3r_\alpha r_\beta - r^2\delta_{\alpha\beta}\bigr)\rho(\mathbf{r})\,d^3r
\end{equation}

\subsubsection{主项：分部积分}

代入 $\rho = \frac{1}{i\omega}\nabla\cdot\mathbf{J}$：
\begin{align}
  Q^e_{\alpha\beta}
  &= \frac{1}{i\omega}\int \bigl(3r_\alpha r_\beta - r^2\delta_{\alpha\beta}\bigr)\nabla\cdot\mathbf{J}\,d^3r \notag\\
  &= -\frac{1}{i\omega}\int \mathbf{J}\cdot\nabla\bigl(3r_\alpha r_\beta - r^2\delta_{\alpha\beta}\bigr)\,d^3r
\end{align}

逐项求梯度：$\nabla(r_\alpha r_\beta) = r_\alpha\hat{\mathbf{e}}_\beta + r_\beta\hat{\mathbf{e}}_\alpha$，
$\nabla r^2 = 2\mathbf{r}$。代入：
\begin{align}
  Q^e_{\alpha\beta}
  &= -\frac{1}{i\omega}\int\Bigl[
    3\bigl(r_\alpha J_\beta + r_\beta J_\alpha\bigr) - 2(\mathbf{r}\cdot\mathbf{J})\delta_{\alpha\beta}\Bigr]\,d^3r
\end{align}

这正是表 1 EQ 的第一项。$\checkmark$

\subsubsection{$k^2$ 修正项（toroidal 四极）}

表 1 EQ 第二项 $\dfrac{k^2}{14}\int\bigl[4r_\alpha r_\beta(\mathbf{r}\cdot\mathbf{J})
- 5r^2(r_\alpha J_\beta + r_\beta J_\alpha) + 2r^2(\mathbf{r}\cdot\mathbf{J})\delta_{\alpha\beta}\bigr]$
是\textbf{环形四极矩}（toroidal quadrupole）的贡献，物理上与电偶极的 toroidal 类似：
描述电流的高阶绕流模式，在 $kr\to0$ 时与电四极辐射不可区分。
（完整推导见 Alaee 2018 补充材料，此处引用其标准结果。）

\subsection{磁四极矩（表 1 MQ）}

磁四极矩的标准定义（对称化）：
\begin{equation}
  Q^m_{\alpha\beta} = \int\Bigl\{
    r_\alpha(\mathbf{r}\times\mathbf{J})_\beta + r_\beta(\mathbf{r}\times\mathbf{J})_\alpha\Bigr\}d^3r
\end{equation}

它来自磁偶极流 $(\mathbf{r}\times\mathbf{J})_\alpha$ 的对称化一阶矩，描述环形磁流模式。

\subsection{表 1 汇总}

\begin{chaptersummary}

\textbf{\S10 章节总结}

\textbf{表 1 的四条公式}：
\begin{itemize}
  \item ED：$p_\alpha = -\frac{1}{i\omega}\int J_\alpha + ikT_\alpha$（含 toroidal 修正）
  \item MD：$m_\alpha = \frac12\int(\mathbf{r}\times\mathbf{J})_\alpha$
  \item EQ：$Q^e_{\alpha\beta} = -\frac{1}{i\omega}\int[3(r_\alpha J_\beta+r_\beta J_\alpha)-2(\mathbf{r}\cdot\mathbf{J})\delta_{\alpha\beta}]$ + toroidal 四极
  \item MQ：$Q^m_{\alpha\beta} = \int\{r_\alpha(\mathbf{r}\times\mathbf{J})_\beta + r_\beta(\mathbf{r}\times\mathbf{J})_\alpha\}$
\end{itemize}

\textbf{推导方法}：\\
(1) 从定义出发（$\mathbf{p}=\int\mathbf{r}\rho$ 等），代入 $\rho = \nabla\cdot\mathbf{J}/(i\omega)$；\\
(2) 分部积分把 $\nabla$ 从 $\mathbf{J}$ 转移到坐标多项式上；\\
(3) toroidal 项（$k^2$ 项）为高阶修正，引用标准结果。

\textbf{局限}：$ka\gtrsim0.5$ 时失效。下一章给出任意 $kr$ 都精确的表 2。

\end{chaptersummary}

\newpage
